\section{Deletion-propagation encoding and a weighted formulation}
\label{app:gdp}

This appendix derives a representation of the temporal problem in the GDP
framework~\citep[Definition 4.1]{makhija2025gdp}, and a weighted ILP extension.
The construction concerns expressibility. It does not transfer the
running-time guarantee of Theorem~\ref{thm:main} to a general-purpose solver.

Use the normalization of Section~\ref{sec:model}. In particular, validate the
original temporal graph before applying forced deletions, set $E_0=E\setminus
F$, and reserve the fixed cost $c(F)$. If $T_0>N_2$, return infeasibility.
Otherwise let $R=N_2-T_0$. The original event identifiers are distinct even
when endpoints and timestamps coincide.

\subsection{Protected compatibility data and unit costs}

Create a unary relation $\mathsf{Event}(i)$ with one row for every identifier
in $E_0$. For $h\in\{2,3\}$, let $H_h$ contain exactly the ordered
$h$-tuples of identifiers forming strictly time-increasing walks ending at
the designated root in $E_0$. These relations have at most $M^2$ and $M^3$
rows. They can be constructed by polynomially many endpoint and
binary-timestamp comparisons; the numeric lifetime is not expanded.

Protect each nonempty $H_h$ by placing its identity query in the GDP
preservation list with threshold $|H_h|$. Since source deletion cannot add
rows to an identity view, no row of that relation can be deleted. An empty
relation needs no protection requirement. The remaining view queries are
ordinary conjunctive queries:
\begin{align*}
 Q_2(i,j)&\mathbin{:\!-}H_2(i,j),\mathsf{Event}(i),\mathsf{Event}(j),\\
 Q_3(i,j,k)&\mathbin{:\!-}H_3(i,j,k),\mathsf{Event}(i),
                 \mathsf{Event}(j),\mathsf{Event}(k).
\end{align*}
Strict comparisons and the root requirement are thus compiled into protected
relation membership, without assuming built-in inequality syntax in GDP.

\begin{proposition}[Unit-cost profile encoding]\label{prop:gdp-encoding}
Suppose all optional events have unit deletion cost. For every
$0\le\ell\le R$ and $0\le w\le N_3$, a GDP instance of size polynomial in
$L$ has a feasible solution exactly when $C(\ell,w)<\infv$; its optimum is
$C(\ell,w)$. An optimal source-deletion set yields an attaining original-event
witness in polynomial time.
\end{proposition}

\begin{proof}
Write $w_2=N_2-\ell$ and $w_3=w$. For each $h\in\{2,3\}$, put $Q_h$ in
the deletion list with requirement $N_h-w_h$ whenever this value is positive,
and in the preservation list with requirement $w_h$ whenever that value is
positive. The former requires at most $w_h$ surviving output rows; the latter
requires at least $w_h$. Omit every zero requirement, consistent with the
positive thresholds in GDP Definition~4.1. At the endpoints, the omitted
inequality is automatic because $0\le |Q_h|\le N_h$.

For each $e\in K$, preserve the Boolean singleton query
$\mathsf{Keep}_e()\mathbin{:\!-}\mathsf{Event}(e)$ with requirement one.
Here $e$ is the constant original identifier. Minimize deletions in the
identity view $Q_E(i)\mathbin{:\!-}\mathsf{Event}(i)$, with no maximization
views. The protected compatibility rows cannot be deleted, and these
singleton requirements prohibit deletion of events in $K$. Therefore every
feasible source-deletion set is exactly a set $D_0\subseteq O$ of optional
event rows.

For any such $D_0$, a row of $Q_h$ survives precisely when every event in its
identifier tuple belongs to $E_0\setminus D_0$. Its cardinality is therefore
$W_h(E_0\setminus D_0)$. All tuple identifiers are output coordinates, so set
semantics does not merge parallel-event walks. A single deleted event affects
all rows containing it in both views. Hence the two exact-count requirements
hold exactly for the deletion sets defining $C(\ell,w)$, and the objective
equals $|D_0|$. There are $O((M+1)^3)$ data rows and $O(M+1)$ view
requirements, with polynomially encoded identifiers and thresholds. This
proves the size, feasibility, optimum, and witness assertions.
\end{proof}

The full deletion witness is $F\cup D_0$ and its total cost is
$c(F)+C(\ell,w)$. The GDP objective above concerns optional deletions only;
forced cost remains the reserved constant.

\subsection{Binary costs and full-profile recovery}

For arbitrary nonnegative binary-encoded costs, use a binary retention
variable $x_e$ for every $e\in E_0$, fixed to one for $e\in K$. Introduce a
binary variable $z_q$ for each tuple $q\in H_2\cup H_3$, and impose
\begin{align}
 z_q&\le x_e &&(q\in H_2\cup H_3,\ e\in q),\label{eq:gdp-and-upper}\\
 z_q&\ge \sum_{e\in q}x_e-|q|+1
              &&(q\in H_2\cup H_3),\label{eq:gdp-and-lower}\\
 \sum_{q\in H_2}z_q&=N_2-\ell,\qquad
 \sum_{q\in H_3}z_q=w.\label{eq:gdp-exact}
\end{align}
The objective is
\begin{equation}
 \min\left(c(F)+\sum_{e\in E_0}c_e(1-x_e)\right).
 \label{eq:gdp-cost}
\end{equation}
For binary variables,~\eqref{eq:gdp-and-upper}--\eqref{eq:gdp-and-lower}
make $z_q$ the conjunction of its event-retention decisions. Thus this ILP
has optimum $c(F)+C(\ell,w)$ whenever the target is attainable, and is
otherwise infeasible. A walk destroyed by several deletions still contributes
only one lost tuple. There are $O((M+1)^3)$ variables and constraints, and
all coefficients have polynomial bit length in $L$. In particular, retaining
$c_e$ as a binary coefficient avoids replication proportional to its value.
This is an explicit weighted witness formulation, not a claim that arbitrary
binary event costs are the literal objective of GDP Definition~4.1.

Solving each target independently and retaining an optimal witness recovers
the full table. The number of target pairs is at most
\[
 (R+1)(N_3+1)\le (M^2+1)(M^3+1)=O((M+1)^5).
\]
Consequently, full-profile output is expressible through polynomially many
exact optimization calls. Their representation size does not bound their
solving time in terms of $R$. The particular tractability results in
\citet[Section 6]{makhija2025gdp} do not automatically transfer under these
joint exact-count requirements and cost coefficients; a proposed transfer
must verify the relevant hypotheses and parameter bounds.

\paragraph{An elementary relaxation diagnostic.}
Let $b_1,b_2:v\to u$ be distinct events at time zero, and let $c:u\to z$ be
an event at time one. All three have unit cost and are optional. This
depth-two instance has $(N_2,N_3)=(2,0)$. With $T_0=1$ and target
$(W_2,W_3)=(1,0)$, the integral optimum is one: delete either middle event;
deleting $c$ misses the target. In the relaxation of the displayed ILP to
$[0,1]$, the assignment
\[
 x_{b_1}=x_{b_2}=1,\qquad x_c=\tfrac12,\qquad
 z_{(b_1,c)}=z_{(b_2,c)}=\tfrac12
\]
is feasible and costs $\tfrac12$. Thus this relaxation is not automatically
integral on the temporal class. This diagnostic concerns the displayed
formulation; it excludes neither stronger formulations nor other algorithms.
It does not replace the structural proof of Theorem~\ref{thm:main} or prove
that no earlier result can imply its guarantee.
